\documentclass[11pt]{article}
\usepackage[margin=1in]{geometry}
\usepackage{amsmath}
\usepackage{booktabs}
\usepackage{microtype}
\usepackage[hidelinks]{hyperref}
\usepackage{parskip}

\title{\vspace{-1.5em}Graph Neural Network Surrogates for a Coupled\\
Heat-Conduction and Pyrolysis Solver in\\ Thermal Protection System Design}
\author{Grey Goodwin \\ \small University of Kentucky}
\date{Draft white paper, 2026-06-15}

\begin{document}
\maketitle

\begin{abstract}
\noindent Thermal protection system (TPS) design relies on repeated evaluation of
an expensive, tightly coupled physics solver that resolves heat conduction,
pyrolysis, pyrolysis-gas transport, and aerothermal surface chemistry. The cost of
those evaluations bottlenecks design-space exploration. We present early work
toward a graph neural network (GNN) surrogate that learns the per-step evolution
of a 1D heat-shield pyrolysis solver, with the goal of running design loops far
faster than the reference code. As a foundation, we re-implement a trusted Fortran
reference solver in modular, vectorized Python and verify it against the original
by column-level comparison of its thermocouple and density outputs. The port
reproduces the reference to machine precision on the conduction-plus-pyrolysis
case and to 0.064\,K on the full gas-plus-aerothermal case. Using data generated
by the validated solver, we then train a per-cell message-passing GNN that
reproduces a held-out heating scenario over a full 60\,s rollout to a mean error
of 7.6\,K, roughly half a percent to one percent of the temperature field, and
tracks the pyrolysis front to within a fraction of a millimeter. This work extends
an existing GNN-surrogate program from canonical flux problems to a multi-physics
engineering solver.
\end{abstract}

\section{Introduction}

Thermal protection systems shield a vehicle from the intense aerothermal heating
of atmospheric entry. Designing one means choosing a material and a thickness that
keep the structure beneath the heat shield below its temperature limit while
adding as little mass as possible. The physics that governs this is a coupled,
strongly nonlinear problem: heat conducts into the material, the material
pyrolyzes (chemically decomposes) as it heats, the pyrolysis gas percolates out
through the charring solid, and the exposed surface exchanges energy with the hot
boundary layer.

A high-fidelity solver that captures this coupling is expensive to run, and TPS
design needs many such runs: parameter sweeps, sizing studies, and optimization
loops each call the solver hundreds or thousands of times. That cost is the
bottleneck this work targets. The response we pursue is a machine-learned
surrogate trained on data from a trusted reference solver, so that design loops
can replace most solver calls with fast surrogate evaluations.

We use a graph neural network for the surrogate because the discretized solver is
naturally a graph: the mesh cells are nodes, and adjacent cells are connected by
edges. A GNN that shares its weights across all cells is mesh-native and, in
principle, applies to a mesh of any size, which is the property we ultimately want
for a tool that must run at varying resolutions.

This paper reports two things. First, a faithful Python re-implementation of the
reference Fortran solver, verified against the original to machine precision on
the simpler test case and to a small fraction of a kelvin on the full
multi-physics case. Second, a first working surrogate trained on data from that
validated solver, accurate over a full rollout on a held-out scenario. We are
explicit about scope: the surrogate is demonstrated here on the
conduction-plus-pyrolysis case, and extending it to the gas-plus-aerothermal case
and to cross-resolution operation are stated as future work, not claimed results.

\section{Reference solver}

The reference is a 1D finite-volume heat-shield code written in Fortran (the
author's \texttt{1Dheat.f} and its associated include files). It models, in one
spatial dimension through the thickness of the slab:

\begin{itemize}
  \item \textbf{Heat conduction} with temperature-dependent and
  char-state-dependent material properties.
  \item \textbf{Pyrolysis}, the heat-driven decomposition of the solid, modeled as
  a set of Arrhenius reactions that consume reactive species and release gas.
  \item \textbf{Pyrolysis-gas transport} through the porous char (Darcy flow plus
  advection), which carries energy and couples back into the temperature field.
  \item \textbf{An aerothermal surface boundary}, including a $B'$
  surface-chemistry table, a blowing correction, and a wall-enthalpy balance,
  representing the interaction with the hot boundary layer.
\end{itemize}

The governing energy equation is the conduction-diffusion balance
\begin{equation}
\rho\, c_p\, \frac{\partial T}{\partial t}
  = \frac{\partial}{\partial x}\!\left(k\, \frac{\partial T}{\partial x}\right)
  + (\text{pyrolysis and gas source terms}),
\end{equation}
discretized on a finite-volume mesh. Each cell carries a temperature and a set of
species densities; the conduction term is evaluated from temperature differences
across cell faces, and the pyrolysis term advances each species by a closed-form
Arrhenius update. Material properties are blended between virgin and char states by
a char-progress variable derived from the local density.

Three reference cases ship with the code and are used throughout: \texttt{aw1}, a
conduction-plus-pyrolysis case with a prescribed wall-temperature history and the
gas physics off, and \texttt{aw2\_tc21} and \texttt{aw2\_tc22}, which add the full
gas and aerothermal boundary. (As shipped, \texttt{aw2\_tc22} is configured
identically to \texttt{aw2\_tc21}; see Section~\ref{sec:val}.)

\section{Methodology}

\subsection{Python port}

The solver was re-implemented in Python one physics module at a time: grid and
geometry setup, mixture properties (conductivity, heat capacity, enthalpy), the
Arrhenius pyrolysis update, the gas solve, and the aerothermal boundary. Where a
coding choice affected the numerical result, the Python deliberately matched the
Fortran. The thermocouple sampling in particular was rewritten to mirror the
reference probe-sampling routine, clamping to the face value at domain boundaries
and reconstructing the face value by interpolation in the interior. Several port
details were matched to specific behavior in the Fortran and materially affected
agreement: thermocouple positions are specified as depth from the right wall; the
face density used for the face conductivity is the pre-pyrolysis density; and
output is written before the physics step, matching the reference time-offset
convention.

\subsection{Verification approach}

Verification is done by direct column-by-column comparison against the Fortran
output file from a run with identical configuration. For every output column
(time, each thermocouple temperature, and each density), we report the maximum
absolute error, the mean absolute error, and the maximum relative error across the
run. This is a strict test: it compares the full time history at every recorded
probe, not just a summary statistic.

\subsection{Training-data generation}

Once verified, the Python solver generates the training data. Each run records the
full per-cell field over time (temperature, bulk density, the per-species
densities, and, in gas runs, the gas pressure and mass flux), including the two
ghost cells at the domain ends that carry the boundary forcing. A single recorded
trajectory can later be turned into training pairs at any time-step spacing without
re-running the solver. The first dataset varies the boundary forcing, a sweep of
wall-temperature histories, on the fast gas-off \texttt{aw1} case.

\subsection{GNN surrogate architecture}

The surrogate is a per-cell message-passing network. Each finite-volume cell is a
node carrying the state vector $[T, \rho, \rho_i, \text{porosity}]$. Adjacent cells
are joined by edges whose features are relative: the difference between the
neighbor's state and the cell's own state, together with the cell spacing. This
relative encoding mirrors the physics, since the conduction flux between two cells
depends on their temperature difference and their spacing.

A forward step encodes the node and edge features with small multilayer
perceptrons, performs one round of message passing (each cell aggregates the
encoded messages from its left and right neighbors and updates its own hidden
state), and decodes a predicted change in the cell's state. The change is added to
the current state and the model is rolled forward in time, exactly as the solver
advances its own update. A single message-passing round is used, matching the
nearest-neighbor reach of the conduction stencil. Because the weights are shared
across all cells, the same trained model applies to a mesh of any length. The
model predicts only the independent degrees of freedom, temperature and the
per-species densities; the bulk density and porosity are derived from the
predicted species densities. Training minimizes the error of a short multi-step
rollout rather than a single step, which controls the accumulation of error over a
long rollout.

\section{Results}

\subsection{Port validation}
\label{sec:val}

\paragraph{aw1 (conduction plus pyrolysis): bit-exact.} All fifteen output columns
match the Fortran output to approximately $10^{-8}$, machine precision, including
every interior thermocouple. The conduction and pyrolysis port is exact.

\paragraph{aw2\_tc21 (full gas plus aerothermal): to 0.064\,K.} With the gas
physics enabled, compared against the same-configuration Fortran output:

\begin{center}
\begin{tabular}{ll}
\toprule
Quantity & Agreement \\
\midrule
All temperature columns & max 0.064\,K, mean about $4\times10^{-4}$\,K \\
Hot aerothermal surface & 0.06\,K \\
Densities & about $10^{-3}$ to $10^{-9}$ \\
Time, cold back face & bit-exact \\
\bottomrule
\end{tabular}
\end{center}

During verification, an apparent early interior drift and a roughly 37.7\,K
back-face offset both proved to be artifacts of how the thermocouples were being
sampled, not numerical error; both vanished once the probe sampling matched the
reference routine. This is a useful caution: comparing two temperature fields
through a sampling layer can manufacture a discrepancy that is not present in the
physics itself.

\paragraph{aw2\_tc22.} As shipped, this case is byte-for-byte identical in
configuration to \texttt{aw2\_tc21} (the same case file, thermocouple locations,
and flux history, and the case-name field still reads \texttt{aw2\_tc21}), so it
does not add an independent validation condition. The Python run reproduces it
identically. Confirmation of whether \texttt{tc22} was intended to carry a distinct
configuration would be valuable.

\subsection{Surrogate accuracy}

Trained on solver-generated trajectories, the surrogate is rolled forward on its
own predictions over a held-out \texttt{aw1} heating scenario (gas off) for the
full 60\,s, with the known boundary forcing re-imposed at each step:

\begin{itemize}
  \item \textbf{Mean rollout error of 7.6\,K}, roughly half a percent to one
  percent of the 298 to 1500\,K temperature field.
  \item The surrogate \textbf{tracks the char and pyrolysis front to within a
  fraction of a millimeter}.
\end{itemize}

Reaching this level required training on multi-step rollouts rather than single
steps. A single-step-trained model is accurate for one prediction but accumulates
error and drifts once it is rolled forward freely; the end-of-trajectory error
fell from several hundred kelvin to tens of kelvin when training was switched to a
multi-step rollout loss. This error-accumulation behavior, and the multi-step
remedy, are consistent with prior surrogate work in this program.

\subsection{Wall time}

The Fortran solver runs the \texttt{aw2\_tc21} case in about 3 minutes; the present
unoptimized Python implementation takes about 15.6 minutes, roughly five times
slower, as expected for interpreted, array-based code matched against compiled
Fortran. The Python port's role is to generate verified training data and to be a
readable reference, not to be fast; the speed argument of this program is the
surrogate, whose wall-time advantage over the solver is the subject of ongoing
measurement.

\section{Discussion}

The results establish two things. First, the multi-physics reference solver can be
reproduced faithfully in a modular, readable Python code: to machine precision on
the conduction-plus-pyrolysis case, and to a small fraction of a kelvin on the full
gas-plus-aerothermal case. That fidelity is what makes the solver usable as a data
generator, since a surrogate is only as trustworthy as the data it learns from.

Second, a per-cell message-passing GNN trained on that data is an accurate
surrogate for the conduction-plus-pyrolysis case over a full rollout, including
faithful tracking of the moving pyrolysis front. The architecture's relative edge
features give it the same local information the physics uses, namely the neighbor
temperature difference and the cell spacing, and its weight sharing across cells is
what makes it mesh-native.

The main limitations are clear and define the path forward. The surrogate is
demonstrated on the gas-off case and on a single mesh resolution. The verification
caution noted above, that a sampling layer can mask or invent discrepancies,
applies equally when comparing a surrogate field against a solver field and is
worth carrying into the surrogate evaluation.

\section{Future work}

\begin{itemize}
  \item \textbf{Gas and aerothermal coverage.} Extend the surrogate to the
  \texttt{aw2} gas-plus-aerothermal case, which adds gas pressure and mass flux to
  the per-cell state.
  \item \textbf{Mesh-resolution generalization.} Train across a range of mesh
  resolutions so that a single model runs at resolutions it did not see in
  training, exploiting the weight sharing the architecture already provides.
  \item \textbf{Speedup measurement.} Quantify the surrogate's wall-time advantage
  over the solver across a design-relevant set of evaluations.
  \item \textbf{Physics extensions.} Add radiation and surface recession, and
  extend the mesh to two and three dimensions.
\end{itemize}

\section*{Acknowledgments}

This work builds on the reference heat-shield solver developed by Dr.\ Martin, and
on an existing graph-neural-network surrogate program within the group.

\end{document}
